% ═══════════════════════════════════════════════════════════════════
% §2  Scalar Types and Type Algebra
% ═══════════════════════════════════════════════════════════════════

\section{Scalar Types and Type Algebra}

\subsection{ScalarType Enumeration}

\begin{definition}[ScalarType]
\label{def:scalartype}
\struct{ScalarType} is an \code{enum class} with underlying type \code{int8\_t}.
Ordinals mirror \code{c10::ScalarType} exactly so that
$\operatorname{static\_cast}\langle\text{int8\_t}\rangle$ is identity between
the standalone library and the PyTorch Vessel adapter.
\end{definition}

\begin{table}[h]
\centering
\begin{tabularx}{\textwidth}{l r r X}
\toprule
\textbf{Name} & \textbf{Ordinal} & \textbf{Bytes} & \textbf{Description} \\
\midrule
\code{Byte}             &  0 & 1  & Unsigned 8-bit integer \\
\code{Char}             &  1 & 1  & Signed 8-bit integer \\
\code{Short}            &  2 & 2  & Signed 16-bit integer \\
\code{Int}              &  3 & 4  & Signed 32-bit integer \\
\code{Long}             &  4 & 8  & Signed 64-bit integer \\
\code{Half}             &  5 & 2  & IEEE 754 half-precision \\
\code{Float}            &  6 & 4  & IEEE 754 single-precision \\
\code{Double}           &  7 & 8  & IEEE 754 double-precision \\
\code{ComplexHalf}      &  8 & 4  & Complex half (re + im) \\
\code{ComplexFloat}     &  9 & 8  & Complex float (re + im) \\
\code{ComplexDouble}    & 10 & 16 & Complex double (re + im) \\
\code{Bool}             & 11 & 1  & Boolean (1 byte storage) \\
\code{BFloat16}         & 15 & 2  & Brain floating-point \\
\code{Float8\_e5m2}     & 23 & 1  & FP8 (5-bit exponent, 2-bit mantissa) \\
\code{Float8\_e4m3fn}   & 24 & 1  & FP8 (4-bit exponent, 3-bit mantissa, no inf/NaN) \\
\code{Float8\_e5m2fnuz} & 25 & 1  & FP8 E5M2 with negative zero as NaN \\
\code{Float8\_e4m3fnuz} & 26 & 1  & FP8 E4M3 with negative zero as NaN \\
\code{Undefined}        & $-1$ & 0 & Sentinel (no type) \\
\bottomrule
\end{tabularx}
\caption{ScalarType enumeration. Ordinals 12--14 and 16--22 are reserved
for future c10 types and must not be used.}
\label{tab:scalartype}
\end{table}

\begin{definition}[\code{element\_size}]
\label{def:element-size}
The function $\operatorname{element\_size} : \text{ScalarType} \to \bv{8}$
is a \code{constexpr} map returning the byte size of each dtype as shown
in the Bytes column of Table~\ref{tab:scalartype}.
For \code{Undefined}, it returns~$0$.
The \code{default} case invokes \code{std::unreachable()}.
\end{definition}

\begin{property}[element\_size totality]
\label{prop:elemsize-total}
$\forall t \in \text{ScalarType} \setminus \{\text{Undefined}\}.\;
\operatorname{element\_size}(t) > 0$. \quad$[\textit{consteval}]$
\end{property}

\subsection{DeviceType Enumeration}

\begin{definition}[DeviceType]
\struct{DeviceType} is an \code{enum class : int8\_t} mirroring \code{c10::DeviceType}.
\end{definition}

\begin{table}[h]
\centering
\begin{tabular}{l r l}
\toprule
\textbf{Name} & \textbf{Ordinal} & \textbf{Description} \\
\midrule
\code{CPU}  &  0 & Host memory \\
\code{CUDA} &  1 & NVIDIA GPU \\
\code{XLA}  &  9 & Google TPU / XLA devices \\
\code{HIP}  & 20 & AMD GPU (ROCm) \\
\bottomrule
\end{tabular}
\caption{DeviceType enumeration (subset relevant to Crucible).}
\end{table}

\subsection{Layout Enumeration}

\begin{definition}[Layout]
\struct{Layout} is an \code{enum class : int8\_t} mirroring \code{c10::Layout}.
\end{definition}

\begin{table}[h]
\centering
\begin{tabular}{l r l}
\toprule
\textbf{Name} & \textbf{Ordinal} & \textbf{Description} \\
\midrule
\code{Strided}   & 0 & Dense strided (default) \\
\code{Sparse}    & 1 & COO sparse \\
\code{SparseCsr} & 2 & CSR format \\
\code{SparseCsc} & 3 & CSC format \\
\code{SparseBsr} & 4 & Block-sparse row \\
\code{SparseBsc} & 5 & Block-sparse column \\
\bottomrule
\end{tabular}
\caption{Layout enumeration.}
\end{table}

\subsection{Strong ID Types}
\label{sec:strongid}

\begin{definition}[\code{CRUCIBLE\_STRONG\_ID}]
\label{def:strongid}
The macro \code{CRUCIBLE\_STRONG\_ID(Name)} generates a struct with the
following exact semantics:

\begin{enumerate}
  \item Storage: a single field \code{uint32\_t v}.
  \item \code{explicit} constructor from \code{uint32\_t} --- no implicit conversion.
  \item Default constructor sets $v = 2^{32}-1$ (\code{UINT32\_MAX}).
  \item \code{static constexpr Name none()} returns the sentinel $2^{32}-1$.
  \item \code{bool is\_valid() const} returns $v \neq 2^{32}-1$.
  \item \code{explicit operator bool() const} delegates to \code{is\_valid()}.
  \item \code{uint32\_t raw() const} returns $v$ (explicit unwrap).
  \item \code{operator<=>} is defaulted (lexicographic on $v$).
  \item No arithmetic operators --- the caller must unwrap, compute, and rewrap.
  \item $\code{sizeof(Name)} = 4$ (verified by \code{static\_assert}).
\end{enumerate}
\end{definition}

Five strong ID types are defined:

\begin{table}[h]
\centering
\begin{tabularx}{\textwidth}{l l X}
\toprule
\textbf{Type} & \textbf{Domain} & \textbf{Semantics} \\
\midrule
\struct{OpIndex}   & $\bv{32}$ & Index into \code{TraceEntry[]} \\
\struct{SlotId}    & $\bv{32}$ & Index into \code{TensorSlot[]} \\
\struct{NodeId}    & $\bv{32}$ & Index into Graph node array \\
\struct{SymbolId}  & $\bv{32}$ & Index into \code{SymbolTable} \\
\struct{MetaIndex} & $\bv{32}$ & Index into \code{MetaLog} buffer \\
\bottomrule
\end{tabularx}
\caption{Strong ID types. All share the same layout and interface
generated by \code{CRUCIBLE\_STRONG\_ID}.}
\end{table}

\begin{invariant}[Sentinel uniqueness]
\label{inv:sentinel}
The sentinel value $2^{32} - 1$ is never a valid index in any array
accessed by these types.  In particular, all Crucible arrays have
fewer than $2^{32} - 1$ elements.
\end{invariant}

\subsection{Strong Hash Types}
\label{sec:stronghash}

\begin{definition}[\code{CRUCIBLE\_STRONG\_HASH}]
\label{def:stronghash}
The macro \code{CRUCIBLE\_STRONG\_HASH(Name)} generates a struct with:

\begin{enumerate}
  \item Storage: a single field \code{uint64\_t v}.
  \item \code{explicit} constructor from \code{uint64\_t}.
  \item Default constructor sets $v = 0$ (``no hash / unset'').
  \item \code{uint64\_t raw() const} returns $v$.
  \item \code{explicit operator bool() const} returns $v \neq 0$.
  \item \code{static constexpr Name sentinel()} returns $2^{64} - 1$ (\code{UINT64\_MAX}).
  \item \code{bool is\_sentinel() const} returns $v = 2^{64}-1$.
  \item \code{operator<=>} is defaulted.
  \item No arithmetic operators.
  \item $\code{sizeof(Name)} = 8$ (verified by \code{static\_assert}).
\end{enumerate}
\end{definition}

Six strong hash types are defined:

\begin{table}[h]
\centering
\begin{tabularx}{\textwidth}{l l X}
\toprule
\textbf{Type} & \textbf{Domain} & \textbf{Semantics} \\
\midrule
\struct{SchemaHash}   & $\bv{64}$ & Op identity (ATen schema) \\
\struct{ShapeHash}    & $\bv{64}$ & Quick hash of input tensor shapes \\
\struct{ScopeHash}    & $\bv{64}$ & Module hierarchy path \\
\struct{CallsiteHash} & $\bv{64}$ & Python source location identity \\
\struct{ContentHash}  & $\bv{64}$ & Region content identity (kernel cache key) \\
\struct{MerkleHash}   & $\bv{64}$ & Subtree identity (includes all descendants) \\
\bottomrule
\end{tabularx}
\caption{Strong hash types.  Default $= 0$ means ``unset''.
Sentinel $= 2^{64}-1$ marks end-of-region boundaries.}
\end{table}

\begin{invariant}[Hash zero and sentinel]
\label{inv:hashzero}
A value of $0$ indicates the hash has not been computed.
A value of $2^{64}-1$ is reserved as a sentinel for boundary marking
and must never be produced by any hash function.
The probability that a uniformly random 64-bit hash equals either
value is $2^{-63}$, which is negligible.
\end{invariant}

All eleven strong types (\struct{OpIndex} through \struct{MerkleHash})
are verified as trivially relocatable on Clang~22 via
\code{CRUCIBLE\_ASSERT\_TRIVIALLY\_RELOCATABLE}, ensuring Arena
\code{memcpy} soundness.

\subsection{Type Algebra: ScalarType Promotion}
\label{sec:type-promotion}

\begin{definition}[Promotion lattice]
\label{def:promotion}
The promotion order $\preceq$ on \struct{ScalarType} is a join-semilattice:
for any two types $a, b$, the \emph{promoted type} is $a \sqcup b = \min\{c : a \preceq c \land b \preceq c\}$.
The lattice satisfies:
\end{definition}

\begin{property}[Promotion properties]
\label{prop:promotion}
For all $a, b, c \in \text{ScalarType}$:
\begin{enumerate}
  \item \textbf{Commutativity:} $a \sqcup b = b \sqcup a$
  \item \textbf{Associativity:} $(a \sqcup b) \sqcup c = a \sqcup (b \sqcup c)$
  \item \textbf{Idempotence:} $a \sqcup a = a$
  \item \textbf{Absorption:} $a \sqcup (a \sqcap b) = a$ where $\sqcap$ is the meet
\end{enumerate}
These properties hold by construction: the promotion function is the
maximum under the partial order induced by the ordinal widening chain
\code{Bool} $\preceq$ \code{Byte} $\preceq$ \code{Short} $\preceq$ \code{Int}
$\preceq$ \code{Long} $\preceq$ \code{Half} $\preceq$ \code{Float}
$\preceq$ \code{Double}, with complex types forming a parallel chain.
$[\textit{consteval}]$
\end{property}
